\documentclass[11pt]{article}
\usepackage[a4paper,margin=25mm]{geometry}
\usepackage{amsmath,amssymb,amsthm,booktabs,hyperref}
\newcommand{\R}{\mathbb R}
\newcommand{\Q}{\mathcal Q_{222}}
\newcommand{\Id}{\mathbb I}
\newtheorem{theorem}{Theorem}
\newtheorem{proposition}[theorem]{Proposition}
\newtheorem{lemma}[theorem]{Lemma}
\title{Explicit semidefinite representations and local certificates for Q222}
\author{}
\date{23 October 2026}
\begin{document}
\maketitle

We use the eight coordinates
\[
 q=(a_0,a_1,b_0,b_1,E_{00},E_{01},E_{10},E_{11}).
\]
The recovered evaluated matrices give an exact representation of a
two-dimensional projection, two explicit local feature constructions,
and four full-coordinate outer approximations. They do not give an
evaluated exact lift of the full eight-coordinate set. This distinction
is independent of the qualitative finite-lift existence theorem.

\section{An exact two-coordinate projection}

Put $t=a_0$ and $s=E_{00}+E_{01}+E_{10}-E_{11}$.
\begin{theorem}
The image of $\Q$ in these coordinates is
\[
 \mathcal T=\{(t,s): |t|\leq1,\quad t^2+s^2/4\leq2\}.
\]
It has the explicit order-four representation
\begin{equation}
 D(t,s)=\operatorname{diag}\left(
 \begin{pmatrix}\sqrt2+t&s/2\\s/2&\sqrt2-t\end{pmatrix},
 1-t,1+t\right)\succeq0.                    \label{eq:direct}
\end{equation}
There are no auxiliary variables. All coefficients belong to
$\mathbb Q(\sqrt2)$, with $\sqrt2$ the positive root of $x^2-2$.
\end{theorem}
\begin{proof}
Binary observables can be dilated to reflections with the same local
embedding, preserving all eight expectations. For reflections define
\begin{align*}
 S&=A_0(B_0+B_1)+A_1(B_0-B_1),\\
 S'&=A_0(B_0-B_1)+A_1(B_0+B_1),\\
 T_\alpha&=\alpha A_0+S,\qquad R_\alpha=\alpha A_1-S'.
\end{align*}
Reflection relations and cross-party commutation give
\[
 T_\alpha^2+R_\alpha^2=(8+2\alpha^2)\Id.
\]
Indeed $S^2+(S')^2=8\Id$, and the linear terms cancel since
$\{A_0,S\}=\{A_1,S'\}$. With $\beta=\sqrt{8+2\alpha^2}$ this gives
\[
 \beta\Id-T_\alpha
 =\frac{(\beta\Id-T_\alpha)^2+R_\alpha^2}{2\beta}\succeq0.
\]
Thus $\alpha t+s\leq\sqrt{8+2\alpha^2}$. For $|t|\leq1$, take
$h=\sqrt{2-t^2}$ and $\alpha=2t/h$. The result is $s\leq2h$.
Flipping both Bob outcomes gives $-s\leq2h$.

Conversely fix $t\in[-1,1]$, put $r=\sqrt{1-t^2}$, and take
\[
 |\psi_t\rangle=\sqrt{\frac{1+t}{2}}|00\rangle+
                  \sqrt{\frac{1-t}{2}}|11\rangle,
 \quad A_0=Z,\quad A_1=X,\quad B_{0,1}=\frac{Z\pm rX}{h}.
\]
These are reflections, and their behavior is
\[
 \left(t,0,\frac t h,\frac t h,\frac1h,\frac1h,
                   \frac{1-t^2}{h},-\frac{1-t^2}{h}\right).
\]
It projects to $(t,2h)$. The Bob flip realizes $(t,-2h)$; a finite
classical mixture realizes every intermediate $s$. At $t=\pm1$ the
same construction remains defined since $h=1$.

Finally \eqref{eq:direct} forces $|t|\leq1$. On this interval its
two-by-two diagonal entries are strictly positive, and its determinant
is $2-t^2-s^2/4$. These conditions prove both matrix inclusions.
\end{proof}

In the original eight coordinates the coefficient matrices of
\eqref{eq:direct} are particularly simple. Write $\mathsf e_{ij}$ for matrix
units with indices $1,\ldots,4$. Then
\[
 \begin{aligned}
 A_0&=\operatorname{diag}(\sqrt2,\sqrt2,1,1),&
 A_{a_0}&=\operatorname{diag}(1,-1,-1,1),\\
 A_{E_{00}}=A_{E_{01}}=A_{E_{10}}
   &=\tfrac12(\mathsf e_{12}+\mathsf e_{21}),&
 A_{E_{11}}&=-\tfrac12(\mathsf e_{12}+\mathsf e_{21}),
 \end{aligned}
\]
and $A_{a_1}=A_{b_0}=A_{b_1}=0$. Their feasible set in $\R^8$ is
the inverse image of $\mathcal T$, not $\Q$.

\subsection{Fixed features and an auxiliary-variable version}

On $\mathcal T$ define the five fixed testing functions
\[
 u_1=\sqrt{\sqrt2+t},\quad u_2=s/(2u_1),\quad
 u_3=\sqrt{\sqrt2-t-u_2^2},\quad
 u_4=\sqrt{1-t},\quad u_5=\sqrt{1+t}.
\]
Their roots are defined on the whole set because
$\sqrt2+t\geq\sqrt2-1>0$ and
$(\sqrt2+t)u_3^2=2-t^2-s^2/4\geq0$. The within-block products in
the spaces $(u_1,u_2)$, $(u_3)$, $(u_4)$, $(u_5)$ are exactly
\[
 u_1^2=\sqrt2+t,\quad u_1u_2=s/2,\quad u_2^2=y,\quad
 u_3^2=\sqrt2-t-y,\quad u_4^2=1-t,\quad u_5^2=1+t.
\]
Thus an extracted lift with $N=5,m=1$ is
\begin{equation}
 L(t,s,y)=\operatorname{diag}\left(
 \begin{pmatrix}\sqrt2+t&s/2\\s/2&y\end{pmatrix},
 \sqrt2-t-y,1-t,1+t\right).                 \label{eq:lift}
\end{equation}
For $(t,s)\in\mathcal T$, the witness
$y=s^2/[4(\sqrt2+t)]$ makes every block positive semidefinite.
Conversely \eqref{eq:lift} gives
$s^2/[4(\sqrt2+t)]\leq y\leq\sqrt2-t$ and $|t|\leq1$,
which imply the defining ellipse inequality. This also proves the
elimination of $y$ without a closure argument.

The same fixed functions certify every affine support. For a source
tilt $\alpha$ and sign $\epsilon=\pm1$, put
\[
 \gamma=\sqrt{4+\alpha^2},\quad
 v_0=\sqrt{(\gamma-\alpha)/2},\quad
 v_1=-\epsilon\sqrt{(\gamma+\alpha)/2}.
\]
Then $v_0v_1=-\epsilon$ and, in separately labelled real closed
source and testing fields,
\[
 \sqrt2\gamma\otimes1-\alpha\otimes t-1\otimes\epsilon s
 =(v_0\otimes u_1+v_1\otimes u_2)^2+(v_1\otimes u_3)^2.
\]
The support value is $h(\alpha)=\sqrt{8+2\alpha^2}$ for
$|\alpha|\leq2$, and $h(\alpha)=2+|\alpha|$ otherwise.
For the latter range, $\sigma=\operatorname{sgn}\alpha$ gives
\[
 2+|\alpha|-\alpha t-\epsilon s
 =4-2\sigma t-\epsilon s+(|\alpha|-2)(1-\sigma t).
\]
The first term uses the displayed square identity at $\alpha=2\sigma$;
the last uses $u_4^2$ or $u_5^2$. Positive constants use
$1=(u_4^2+u_5^2)/2$. Scaling these slacks and adding a positive
constant covers every affine function nonnegative on $\mathcal T$.
This includes unbounded source tilts and all tilt endpoints without
identifying the two field copies.

\subsection{An exact obstruction to treating this as the full set}

The rational behavior
\[
 q_*=(0,0,0,0,1,1,1/4,-1/4)
\]
has nonnegative nonsignalling probabilities
$p(a,b|x,y)=(1+abE_{xy})/4$. Its maximum absolute CHSH value over
the eight sign patterns is $5/2<2\sqrt2$. Every marginal is zero,
so it satisfies every standard relabelled single-marginal tilted-CHSH
inequality, whose bound is at least $2\sqrt2$. It also lies in
\eqref{eq:direct}, since $t=0,s=5/2$ and $s^2/4=25/16<2$.

It is not quantum. After dilation and purification, $E_{00}=E_{01}=1$
would imply
\[
 \|(A_0-B_y)\psi\|^2=2-2E_{0y}=0\quad(y=0,1).
\]
Hence $(B_0-B_1)\psi=0$, forcing $E_{10}-E_{11}=0$, whereas the
displayed point gives $1/2$. Thus intersecting all these projected
lifts and the probability inequalities still does not recover $\Q$.

\section{An evaluated regular-C2 local certificate}

The following construction recovers a complete local coefficient lift.
Its domain is part of the assertion. For the seven-source order
$(d,n,m,C,A,D,\alpha)$, the full quarter-parent is
\begin{align*}
 \widetilde P={}&nD_uD_vX^2+m(vW-X+uvY-uZ)^2+CD_v(uW-Y)^2+AD_uD_vY^2\\
 &+(D+d/2)(uW+uvX-Y-vZ)^2+(\alpha+d/2)D_u(vY-Z)^2\\
 &+(d/2)D_v(X+uZ)^2-(d/2)D_uD_vZ^2,
 \qquad D_u=1+u^2,\quad D_v=1+v^2.
\end{align*}
Set only $d=Z=1$, and write $P_c(z)$, where
$c=(n,m,C,A,D,\alpha)$ and $z=(u,v,W,X,Y)$. The exact anchors are
\[
 c_*=(7/8,9/8,1/18,11/8,9/8,1/18),\qquad
 \theta_*=(1,1,1,-1/10,1/10).
\]
Let $F(c,\theta)=(P_c,\partial_uP_c,\partial_vP_c,
\partial_WP_c,\partial_XP_c,\partial_YP_c)(\theta)$.
Put $s=(c-c_*,h)$, $h=\theta-\theta_*$, and independently
$t=z-\theta_*$. The source domain and testing box are exactly
\[
 F(c,\theta)=0,\quad |s_i|\leq\rho\ (i=1,\ldots,11),\qquad
 |t_j|\leq\rho\ (j=1,\ldots,5),\qquad
 \rho=\frac{510101005}{350806159660744}.
\]
The value equation supplies the remaining homogeneous $Z$ derivative
by Euler's identity. The six-by-six source Jacobian has determinant
$114048/78125$ at the anchor. On this entire contact box its relative
row-norm perturbation is at most
$11650604933999/11225797109143808<1$, so the source solution is unique
for any fixed contact in the box. No global cover is implied.

Here is an exact finite prescription for all the features and matrices.
Expand $P_c(\theta+\Delta)$ and remove degree zero and one in
$\Delta$. Allocate each remaining monomial to a symmetric matrix $H$:
use the first coordinate with exponent at least two as a diagonal
square; if none exists, use the first two occupied coordinates and
put half the coefficient in each symmetric off-diagonal entry. Then
\[
 P_c(\theta+\Delta)=P_c(\theta)+\nabla P_c(\theta)^T\Delta+
                                      \Delta^TH\Delta
\]
is an exact polynomial identity with no Taylor remainder. At the anchor
$H_0=LD_0L^T$, where
\[
 L=\begin{pmatrix}
 1&0&0&0&0\\
 -9801/10201&1&0&0&0\\
 1210/10201&8991/20002&1&0&0\\
 10080/10201&-5040/10001&79/1010&1&0\\
 -10120/10201&5060/10001&-119/1010&2/81&1
 \end{pmatrix},
\]
\[
 D_0=\operatorname{diag}\left(\frac{10201}{3600},
 \frac{20002}{91809},\frac{27775}{10001},\frac{891}{202},\frac{1067}{162}\right).
\]
Thus $f=L^Tz$, $\widehat H=L^{-1}HL^{-T}$ give the correctly oriented
identity $\Delta^TH\Delta=(\Delta f)^T\widehat H\Delta f$.
Put $\varepsilon=10001/918090$. Use the 25 linear forms
$\ell_k=f_i$ and $f_i\pm f_j$, in that order, with coefficients
\[
 \widehat H_{ii}-8\varepsilon,\qquad
 \varepsilon+\widehat H_{ij}/2,\qquad
 \varepsilon-\widehat H_{ij}/2.
\]
Each coordinate participates in four pairs, so the pair diagonal
contribution is exactly $8\varepsilon$. Substitute $\Delta=t-h$ and
write the resulting coefficient polynomials as $\gamma_k(s,t)$.
Their constants are at least $\varepsilon$, and their largest
nonconstant coefficient norm is exactly
\[
 B_*=\frac{43803942777143}{46827180450},\qquad
 \rho=\frac{\varepsilon}{8(1+B_*)}<\tfrac12.
\]
All these rational coefficient data are given in the accompanying
\path{regular-c2-local/regular_c2_frames.json}.

For completeness, the separation of every coefficient into nonnegative
pure factors has a uniform exact formula. Group its support uniquely as
\[
 \gamma=\gamma_0+G(s)+T(t)+\sum_{\beta\ne0}A_\beta(s)t^\beta,
\]
where $G,T,A_\beta$ have zero constant. Define
\[
 \sigma_\beta=\sum_a|[s^a]A_\beta|,\quad
 U_\beta=A_\beta+\sigma_\beta\rho,\quad
 V_\beta=t^\beta+\rho^{|\beta|},\quad
 B=\sum_\beta\sigma_\beta\rho^{|\beta|+1},
\]
\[
 \begin{aligned}
 p(s)&=\gamma_0/2-B/2+G(s)-\sum_\beta\rho^{|\beta|}A_\beta(s),\\
 q(t)&=\gamma_0/2-B/2+T(t)-\sum_\beta\sigma_\beta\rho\,t^\beta.
 \end{aligned}
\]
Expansion proves $\gamma=p+q+\sum U_\beta V_\beta$.
On the independent boxes, $U_\beta,V_\beta\geq0$. If $g,r,M$ are
the coefficient norms of $G,T,\sum A_\beta t^\beta$, then
$g+r+M\leq B_*$, $B\leq\rho M$, and
\[
 p,q\geq\gamma_0/2-\tfrac32\rho(g+r+M)
       \geq5\varepsilon/16>0.
\]
These inequalities prove all root signs, including box boundaries.

For each frame and each pair $(U,V)=(p,1),(1,q),(U_\beta,V_\beta)$,
use the testing vector and source matrix
\[
 \psi_b(t)=\sqrt{V(t)}\binom{1}{\ell_k(t)},\qquad
 Q_b(s)=U(s)\begin{pmatrix}\ell_k(h)^2&-\ell_k(h)\\-\ell_k(h)&1\end{pmatrix}.
\]
Then $Q_b\succeq0$ and
$\psi_b^TQ_b\psi_b=UV(\ell_k(t)-\ell_k(h))^2$. The 735 mixed
pairs and 50 pure pairs give exactly 785 blocks of order two and
1570 listed testing functions. Contact makes the discarded constant
and linear Taylor terms zero, proving the full-parent identity.
The same proof is a tensor sum of squares in independently labelled
real closed fields because each square root is taken in its own field.

The within-block products are $V,V\ell_k,V\ell_k^2$. Let
$m_0,\ldots,m_{236}$ be the lexicographically ordered union of their
monomial supports. Define exact rational arrays by coefficient extraction
\[
 \psi_{b,i}\psi_{b,j}=\sum_kp_{b,ij,k}m_k,
 \qquad \widetilde P_{\lambda}(\theta_*+t,1)
             =\sum_{k,\nu}T_{k\nu}\lambda_\nu m_k,
 \quad\lambda=(d,n,m,C,A,D,\alpha).
\]
The evaluated data contain 27747 nonzero product coefficients and
295 nonzero target coefficients. The exact coefficient lift is
\[
 \sum_{b,i\leq j}(2-\delta_{ij})p_{b,ij,k}G_{b,ij}
       =\sum_\nu T_{k\nu}\lambda_\nu\quad(k=0,\ldots,236),
 \qquad G_b\succeq0.
\]
It has 2355 upper-triangular Gram variables and 237 equations.
Every feasible point gives the displayed nonnegative polynomial
identity; every source in the contact box has the explicit feasible
matrices above. Equality with the entire cone of polynomials nonnegative
on the box is not asserted. This is a local certificate lift in source
coordinates, not an eight-coordinate quantum-set lift.

\section{A second recovered local feature library}

The September C1D construction gives a different explicit finite local
library. Its source domain is
\[
 0<a\leq1,\quad0<b,q\leq\epsilon,\quad0<r<1,\quad0<E\leq1,
 \qquad\epsilon=1/2788018,
\]
with independent tests $|u|,|v|,|W|,|X|,|Y|\leq\epsilon$, $Z=1$.
Put $e^2=E$, $r_r^2=r$, $s_s^2=1-r$ with positive roots. Its five
separated roots are
\[
 \begin{pmatrix}
 uW-b^2v\\buv\\b^2(v^2-q^2)\\
 (b/e)(X+u)-a^2b^3q^2e\,vW\\(Y-ruW)/(r_rs_s)
 \end{pmatrix}.
\]
The recovered full source Gram is $8I_5$ plus a polynomial correction
whose upper coefficient norm is $1394008$. Each correction monomial
contains a small source or testing variable. Hence its row perturbation
is less than $1/2$ on this rectangle, uniformly without lower bounds on
$E,r,1-r$.

For a scalar channel $\rho+\sum\beta_mC_m\otimes T_m$, set
$S=\sum|\beta_m|$. Choose the small pure factor as $t_m$, absorb
the sign into it, and call the other factor $l_m$. Each term has
$|t_m|\leq\epsilon$, $|l_m|\leq1$, and $\rho/S>\epsilon$.
The exact identity
\[
 \rho/S+t_m\otimes l_m
 =\tfrac12\{(\rho/S+t_m)\otimes(1+l_m)
              +(\rho/S-t_m)\otimes(1-l_m)\}
\]
uses nonnegative factors on their own sides. The 25 diagonal and pair
channels produce 82928 tensor squares. The data specify all 41464 term
pairs and all pure-factor radical choices.

After adjoining the 1212 testing radicals, the fixed testing spanning
functions are the 3844 listed pairs of
\[
 \kappa\zeta_j\chi_a,\qquad
 \chi_a\in\{1,uW,v,uv,v^2,X+u,vW,Y\},
\]
where
$\kappa^2(1+u^2)(1+v^2)(1+W^2+X^2+Y^2)=1$.
Box equations $\operatorname{box}_t^2+t^2=\epsilon^2$ and the
printed radical equations give 1223 algebra variables and 1218 equations.
The testing image is this entire closed local box. The list is not
claimed linearly independent. Neither its union with the regular-C2
library nor its own local identities establishes a finite global cover.
The original polynomial allocation multipliers and pure-factor
restoration are preserved beside the library; they are required to
identify its source Gram with the original Bell slack.

\section{Evaluated full-coordinate outer pencils}

The recovered data include the following actual rational matrices.
Here $N$ is the total block-diagonal order and $m$ is the auxiliary
dimension; none of these rows is an evaluated exact full-set lift.
\begin{center}
\begin{tabular}{rrrl}
\toprule
$N$&$m$&nonzero records&reverse inclusion\\\midrule
8368&16480&55712&disproved\\
10736&21472&72608&disproved\\
13744&29920&153248&unproved\\
63888&31111&455750&unproved (C160)\\\bottomrule
\end{tabular}
\end{center}
For the first three files a record $[k,i,j,a]$ adds $a$ to the upper
entry $(i,j)$ of $A_0$ if $k=0$, of $A_k$ if $1\leq k\leq8$, and
of $B_{k-9}$ otherwise. Mirror each off-diagonal once. The represented
set is
\[
 \mathcal O_N=\{q:\exists y\in\R^m,
                A_0+\sum_{k=1}^8q_kA_k+\sum_{j=0}^{m-1}y_jB_j\succeq0\}.
\]
Here the auxiliary variables and matrices are indexed from zero,
so a coefficient index $k=9+j$ multiplies $y_jB_j$.
C160 instead stores $[b,i,j,k,p,q]$, meaning the entry $p/q$ in
block $b$ and coefficient $k$ with the same visible/auxiliary convention.
Its block orders and all 455750 records are supplied explicitly.

The first three constructions dualize finite cones of globally valid
source certificates, with exact source maps, all source signs, and
explicitly positive strict source points. An elementary adjoint argument
then proves $\Q\subseteq\mathcal O_N$: if $L(c,z)$ is a source
certificate pencil with a strict positive point, the image
$\{L^*(Y):Y\succeq0\}$ is closed. Indeed bounded adjoint images,
paired with the strict source point, bound $\operatorname{tr}Y$;
subsequential limits give closedness. Bipolarity identifies the dual of
the projected source cone with the attained adjoint image. Exact linear
elimination supplies every affine auxiliary coordinate in the recovered
primal arrays. Thus source-certificate soundness gives the quantum
forward inclusion, not its converse.

For $N=13744$ the source certificates use four blocks of orders
$(14,15,14,15)$ for C1 and $(15,14,14,15)$ for C2 in the full identity
\[
 (1+u)(1+v)F_c=\sum_{g\in\{0,1\}^2}u^{g_1}v^{g_2}
                                       f_g^TG_gf_g,\quad G_g\succeq0.
\]
The multiplier is pure testing and positive for $u,v>0$; state sign
gauges and homogeneous limits give the other charts. The global pencil
retains 128 C1 copies, 64 C2 copies, 64 adjacent and 32 opposite exact
closed-face copies, and 16 event inequalities. Their local orders and
auxiliaries give exactly $N=13744,m=29920$, with 2928 blocks of maximum
order 15. Adequacy of these fixed features for every physical source is
not established by their coefficient matching identities.

\subsection{Why the two smaller global candidates are strictly outer}

The exact physical C1 source is
\[
 \begin{aligned}
 d&=1,&m&=\frac{1166523000}{31661820415017653},&
 C&=\frac{322803000}{186199784281},\\
 A&=\frac{218288946200}{1359122513},&
 D&=\frac{90943434293}{222507771414},&
 \alpha&=\frac{248802430236}{3300726103},\\
 \beta&=\frac{205476973}{51646655494}.&&&&
 \end{aligned}
\]
It has the contact $(u,v;\psi)=(1/13,16;140,-1/32,1/50,1)$.
Its complete state-kernel and angular equations hold exactly. Its six
positive Laurent determinant weights balance around zero, so the
elementary weighted AM--GM identity proves determinant nonnegativity.
The positive $m,C,A$ rows have a strict complementary state block;
the Schur complement therefore proves physicality, including signs
and homogeneous endpoints.

Let $f_*=J_1c_*$, where
\[
 \begin{aligned}
 (J_1c)_0&=d+m+C+A+D+\alpha+\beta,\\
 (J_1c)_{1:4}&=(-A-\alpha,m-C-D+\beta,C+A-\beta,-m+D-\alpha),\\
 (J_1c)_{5:8}&=(-d/2-A,d/2+\alpha,-d/2-C-\beta,-d/2-m-D).
 \end{aligned}
\]
The four kernel equations, two angular equations and two deterministic
zeros have rank eight in the nine affine coefficients: the indicated
eight-column minor is
\[
 \frac{409072768794400177096031164572864}{10770194675703125}\ne0.
\]
In any positive physical decomposition every summand obeys these same
eight equations. It is therefore proportional to $f_*$, with a
nonnegative scalar. Thus $f_*$ is an extreme physical ray.

For the actual order-$8368$ and order-$10736$ source cones the only
nonnegative component preimages of this ray are $c_*$ and
$\sigma_1c_*=(d,C,m,\alpha,\beta,A,D)$. The complete source-map
arrays retain the finite screen establishing this statement. The
preserved positive rational dual matrices give, on the actual C1
coefficient pencil, the trace identities
\[
 \sum_b\operatorname{tr}(Y_b^{(0)}L_b(c,y))=\ell(c)-d,\qquad
 \sum_b\operatorname{tr}(Y_b^{(1)}L_b(c,y))=\ell(\sigma_1c)-d,
\]
where $\ell(c)=93424\,\psi_*^TP_c(u_*,v_*)\psi_*-d$ and
$\ell(c_*)=-1$. Every auxiliary coefficient is zero, and all dual
matrices are positive definite. The Gram trace at the proposed source
is therefore $-2$: the further subtraction of $d$ is the actual scalar
certificate coefficient, not a change of parent normalization. Hence
neither permitted preimage lies in its certificate cone.

These finite source sums are closed: positive finite evaluation moments
bound all Gram matrices of convergent certificate sources, and the
positive affine constant bounds each nonnegative component in the sum.
Thus separating $f_*$ from the closed source sum gives a point of its
dual violating $f_*$. The event rays force its constant coordinate to
be positive, so it normalizes to $(1,q)$ with $q\notin\Q$.
Consequently
\[
 \Q\subsetneq\mathcal O_{8368},\qquad
 \Q\subsetneq\mathcal O_{10736}.
\]
This proof gives strictness; it does not print one numerical behavior
and thousands of auxiliary coordinates. It does not refute the enlarged
order-$13744$ or the separate C160 candidate.

\subsection{The C160 moment construction and its coverage dependency}

For the fourth candidate use angular parameters
$z=(k,x_0,y_0,x_1,y_1)\in[0,1]^5$ with
\[
 ky_0(1-x_0)\geq x_0(1-y_0),\quad
 kx_1(1-y_0)\geq y_0(1-x_1),\quad
 ky_1(1-x_1)\geq x_1(1-y_1).
\]
Call this domain $T$. Put $d(t)=(1-t)^2+t^2$,
$c(t)=(1-2t)/d(t)$, $s(t)=2t(1-t)/d(t)$ and
$D=(1+k^2)d(x_0)d(y_0)d(x_1)d(y_1)$, so $1/16\leq D\leq2$.
The physical map is
\[
 \Phi_{A_x}=\frac{1-k^2}{1+k^2}c(x_x),\quad
 \Phi_{B_y}=\frac{1-k^2}{1+k^2}c(y_y),\quad
 \Phi_{E_{xy}}=c(x_x)c(y_y)+\frac{2k}{1+k^2}s(x_x)s(y_y).
\]
Its stated full-quantum coverage is
\begin{equation}
 \Q=\operatorname{conv}\bigl(\{\ell_1,\ldots,\ell_{16}\}
            \cup\bigcup_{g\in\mathcal G}g\Phi(T)\bigr),
 \label{eq:c160-coverage}
\end{equation}
where $\mathcal G$ contains the 128 local relabellings and party swaps.
The recovered construction invokes the Barizien--Bancal extremal
characterization~\cite{BB} for \eqref{eq:c160-coverage}. This external coverage
theorem and its coordinate interface are a dependency of its full-$\Q$
forward inclusion; they were not independently re-proved in this
extraction. The following moment statement itself is direct.

Let $p_e=\prod_i p_{e_i}(z_i)$ for
$(p_0,p_1,p_2)=((1-t)^2,t(1-t),t^2)$ and $e\in\{0,1,2\}^5$.
There are 243 moments. The 160 printed feature blocks have entries
$\sqrt{w_b/D}$ times products of $z_i$ or $1-z_i$, where every
$w_b$ is a nonnegative rational multiple of the cube factors and the
three displayed domain inequalities. Their orders sum to 499.
The source cache prints all product coefficients in the $p_e$ basis
and all nine target columns $D,-D\Phi_1,\ldots,-D\Phi_8$.
For $z\in T$, the exact witness $Y_e=p_e(z)/D(z)$ makes every moment
block the rank-one Gram of its feature vector, and the target map gives
$(1,\Phi(z))$. Nine independent target equations eliminate nine moments,
leaving 234 canonical auxiliaries. Thus
$\operatorname{conv}\Phi(T)$ lies in the canonical order-499 projection.

Block zero has weight one, includes all five variables, and its products
span all 243 moments. The exact constant-column Gram
\[
 Q_D=\begin{pmatrix}I_{16}&I_{16}\\I_{16}&2I_{16}\end{pmatrix}\succ0
\]
bounds the trace of its PSD moment block at mass one and forces every
moment to vanish at mass zero. Therefore the canonical projection is
compact and its homogenization introduces no nonzero zero-mass point.
Taking 128 copies and 16 nonnegative deterministic weights, then
eliminating nine independent weights, gives
$N=128\cdot499+16=63888$ and $m=128\cdot243+7=31111$.
This proves the forward inclusion subject to
\eqref{eq:c160-coverage}. The reverse inclusion would require a
universal adequacy theorem for this particular fixed 160-block space;
neither the Dirac witnesses nor exact elimination supplies it.

\section{The effective construction that remains unevaluated}

There is a precise finite search, conditional only on the independently
established qualitative existence theorem. The full quantum set has a
rational first-order membership formula using nine real two-qubit atoms:
each atom has a unit real four-state vector and two unit measurement
circles, and its eight outputs are Pauli expectation polynomials. Nine
convex weights suffice by Carath\'eodory in $\R^8$. This gives 81
existential variables, 36 equations, nine inequalities and degree at most
five, covering complex mixed models by realification and direct sums.

For a proposed single PSD order $N$, set $D=N(N+1)/2$. Any auxiliary
matrix span compresses to at most $D$ variables. Let $c$ contain all
unknown entries of $A_0,A_1,\ldots,A_8,B_1,\ldots,B_D$. The sentence
\[
 \exists c\ \forall q\quad
 [\operatorname{Phys}(q)\ \longleftrightarrow\
       \exists y\in\R^D:L_c(q,y)\succeq0]
\]
is a rational first-order sentence, with PSD replaced by all principal
minors. Quantifier elimination decides it. If true over $\R$, it is true
over the real algebraic numbers, so exact algebraic coefficient witnesses
can be extracted with real-root selection data. Enumerating $N=1,2,\ldots$
terminates if some finite lift exists. This is an effective procedure;
it supplies no evaluated order, coefficient tuple or useful running-time
bound until executed.

The stored three-block order-two instance has 162 unknown coefficients
and nine auxiliaries and is explicitly unsolved. Existence of an arbitrary
finite SDP lift does not ensure termination of an order-two-only search.

The unevaluated global tasks are therefore concrete: E1 requires the
selected global testing functions, product basis and dimensions; E2
requires one fixed exact original-coordinate pencil; E3 requires both
full-set inclusions for that same pencil. Local certificates, rejected
outer candidates and an executable first-order input do not by themselves
complete these tasks.

\section*{Scientific data and exact checks}

The accompanying \path{extracted-lifts/} directory contains the exact
projection matrices and checker; the complete regular-C2 feature and
product arrays with an independent standard-library checker; the C1D
local square library, source multipliers and algebra; and all four global
outer matrices with source maps and separator data. Its README specifies
the two COO conventions and the status of each construction. In the
two disproved candidate files only scientific scope metadata was updated;
all coefficient and block arrays are unchanged. Provenance records bind
the source bytes and the copied mathematical data. No matrix here is
presented as an exact full-$\Q$ lift.
\begin{thebibliography}{1}
\bibitem{BB} V. Barizien and J.-D. Bancal,
\emph{Quantum statistics in the minimal scenario},
\href{https://arxiv.org/abs/2406.09350}{arXiv:2406.09350} (2024),
Theorem 1 and its full-alternance parametrization.
\end{thebibliography}
\end{document}
